\documentclass[../../main.tex]{subfiles}

\begin{document}

\chapterstarstyle{Abstract}{\textsc{Preamble}}
\vspace{1.4cm}

Superconductivity is a paradigmatic example of a many-body phenomenon, where collective electronic correlations give rise to macroscopic quantum coherence. To capture these correlations beyond mean-field and purely local descriptions, diagrammatic extensions of dynamical mean-field theory (DMFT) have been developed. Among them, the dynamical vertex approximation ($\dga$) incorporates non-local correlations by using the local two-particle irreducible vertex as a building block, thereby providing access to momentum- and frequency-dependent self-energies, susceptibilities and more. This makes the dynamical vertex approximation particularly well-suited to describe phase transitions such as superconductivity, where non-local fluctuations are essential.
\newline

In this thesis, we present a code for the self-consistent ladder dynamical vertex approximation applied to the multi-orbital Hubbard model. The implementation exploits vertex asymptotics, enabling calculations down to very low temperatures and thus the determination of superconducting transition temperatures from two-particle correlations and their feedback into single-particle properties. Compared to DMFT, which neglects all spatial correlations, this framework systematically incorporates them and provides a natural extension to the local theory.
\newline

A central focus of this work is the multi-orbital nature of correlated superconductors. While single-orbital models capture only basic aspects, realistic materials often exhibit orbital degrees of freedom that crucially shape pairing mechanisms. For example, the recently discovered infinite-layer nickelates (e.g. NdNiO$_2$) display significant multi-orbital character that cannot be described sufficiently within a single-orbital framework. Our numeric implementation of a multi-orbital dynamical vertex approximation can help capturing orbital selectivity, Hund's coupling and intertwined instabilities, which are decisive, among other factors, for the superconducting pairing structure and critical temperature.
\newline

By combining methodological advances with a focus on orbital complexity, this thesis establishes a computational framework for studying unconventional superconductivity in multi-orbital systems, with direct relevance for cuprates, nickelates, iron-based superconductors and other strongly correlated materials.

\end{document}